% \iffalse meta-comment
%
% hit-final-paragraph.dtx 是段落与行距
%
% \fi
%
% \section{段落与行距}
%
% \changes{v3.2a}{2026/09/28}{列表照 iota：编号顶在首行缩进处、续行回左缘，中文“（1）”后面不空，英文版与博士后
%   “(1)”加一个空格；先前编号右对齐在 3\,em／3.5\,em 处}
% \changes{v3.2a}{2026/09/13}{列表间距迁进 \file{src/engine/hit-list.dtx}，这里只留正文首行
%   缩进，量的从西文 em 换成汉字宽 \texttt{2\cs{ccwd}}}
%
% ^^A ======================================================================
% ^^A Module final-paragraph: 段落首行缩进与标点处理（hit-thesis）
% ^^A ======================================================================
%    \begin{macrocode}
%<*final-paragraph>
\bool_if:NT \g__hit_final_bool
  {
\ctexset{%
  punct=quanjiao,
  space=auto,
  autoindent=true}
%    \end{macrocode}
%
% 列表照 iota：编号顶在首行缩进处、续行回到版心左缘，与范例里手敲的“（1）……”段落一个样
% （两份范例的 docx 里一个 \texttt{w:numPr} 都没有）。中文是全角“（1）”，后面不空；英文版与博士后
% 是半角“(1)”加一个空格。机制在 \file{src/engine/hit-list.dtx}（Word 的多级列表，\pkg{enumitem} 照它排）。
% 先前是编号右对齐在 3\,em（列表）、3.5\,em（编号列表）处，后面空半个 em，没有来路。
%    \begin{macrocode}
\bool_lazy_or:nnTF { \g__hit_en_bool } { \g__hit_postdoc_bool }
  { \hit@list@word:nn { ( \% 1 ) } { space } }
  { \hit@list@word:nn { （ \% 1 ） } { nothing } }
%    \end{macrocode}
%
% 正文首行缩进。规范要的是「空两个汉字」，所以量 \cs{ccwd}（\pkg{ctex} 的汉字宽）
% 而不是 \texttt{em}（当前\emph{西文}字体的 em）。两个数在这里恰好相等，
% 宋体一个字与 Times 的一个 em 都等于字号，所以一直没出岔子。
%
% 但拴在西文字体上会出事：西文字体加载失败时退回 \cs{nullfont}，它的
% \cs{fontdimen}~6 是零，\texttt{2em} 静默变成 0\,pt，正文首行缩进整个消失，
% 日志里只有一句 \texttt{not loadable} 混在一堆 \pkg{fontspec} 报错里。
% \pkg{newtx} 没装（\option{fontset} 取 \texttt{fandol}、\texttt{ubuntu} 时
% 西文走 \texttt{TeXGyreTermesX}，那个字体是 \pkg{newtx} 带的）就会这样。
% report 类那边 \issue[201] 已经改成 \texttt{2\cs{ccwd}} 了，这里跟上。
%
% 新版面不走这一句：\cs{docxlike\_format\_apply\_body:} 会用字符网格的步进重设
% \cs{parindent}，那才是 Word 的算法（范例存的是
% \texttt{firstLineChars="200"} 加 \texttt{firstLine="498"} 缇，
% 498/20/2 = 12.45\,bp 一个字符，等于字号加字距，不是字号）。
%    \begin{macrocode}
\dim_set:Nn \parindent {2\ccwd}
  }
%</final-paragraph>
%    \end{macrocode}
%
% \Finale
